% ============================================================================
%  Packages, theorem environments and presentation styling.
%
%  Mathematical *notation* lives in notation.tex, not here. The split is
%  deliberate: change how the document looks in this file, change what the
%  symbols are in that one.
% ============================================================================

% ---------------------------------------------------------------------------
%  Encoding, fonts, page geometry
% ---------------------------------------------------------------------------
\usepackage[T1]{fontenc}
\usepackage[utf8]{inputenc}
\usepackage{lmodern}
\usepackage{microtype}

\usepackage[a4paper,margin=3cm,marginparwidth=2cm]{geometry}

% ---------------------------------------------------------------------------
%  Mathematics
% ---------------------------------------------------------------------------
\usepackage{amsmath}
\usepackage{amssymb}
\usepackage{amsthm}
\usepackage{mathtools}
\usepackage{stmaryrd}     % \llbracket .. \rrbracket for semantic brackets
\usepackage{mathpartir}   % \inferrule and the mathpar environment

% ---------------------------------------------------------------------------
%  Tables, lists, graphics
% ---------------------------------------------------------------------------
\usepackage{array}
\usepackage{booktabs}
\usepackage{longtable}
\usepackage{tabularx}
\usepackage{multirow}
\usepackage{enumitem}
\usepackage{xcolor}
\usepackage{graphicx}

\usepackage{tikz}
\usetikzlibrary{positioning,arrows.meta,calc,fit,backgrounds,%
                shapes.geometric,matrix,decorations.pathreplacing,patterns}

% ---------------------------------------------------------------------------
%  Colour palette
% ---------------------------------------------------------------------------
\definecolor{dstmblue}{HTML}{1F4E79}
\definecolor{dstmteal}{HTML}{116B70}
\definecolor{dstmred}{HTML}{9B2226}
\definecolor{dstmamber}{HTML}{9C6500}
\definecolor{dstmgrey}{HTML}{4A4A4A}
\definecolor{resbg}{HTML}{EEF4FA}
\definecolor{devbg}{HTML}{FDF1F1}
\definecolor{codebg}{HTML}{F6F6F4}

% ---------------------------------------------------------------------------
%  Source code listings
% ---------------------------------------------------------------------------
\usepackage{listings}
\lstdefinestyle{dstmc}{%
  language=C,
  basicstyle=\ttfamily\footnotesize,
  keywordstyle=\color{dstmblue}\bfseries,
  commentstyle=\color{dstmgrey}\itshape,
  stringstyle=\color{dstmteal},
  backgroundcolor=\color{codebg},
  frame=leftline,
  framerule=1pt,
  rulecolor=\color{dstmblue!30},
  framesep=6pt,
  xleftmargin=8pt,
  breaklines=true,
  breakatwhitespace=true,
  columns=fullflexible,
  keepspaces=true,
  showstringspaces=false,
  morekeywords={dstm_time_t,dstm_ptr_t,dstm_darray_t,dstm_dgraph_t,%
                dstm_distance_t,dstm_region_t,dstm_opts_t,size_t,%
                int64_t,intptr_t,MPI_Comm},
}
\lstset{style=dstmc}

% ---------------------------------------------------------------------------
%  Callout boxes
% ---------------------------------------------------------------------------
\usepackage[most]{tcolorbox}
\tcbuselibrary{skins,breakable}

% ---------------------------------------------------------------------------
%  Theorem environments
%
%  Two numbering families:
%
%    * Mathematical results share one per-chapter counter, so a cross
%      reference such as "Theorem 2.7" is unambiguous without the reader
%      having to know which of five counters it came from.
%
%    * Restrictions (O1..) and deviations (D1..) get flat document-wide
%      counters with a letter prefix, because they are referred to by name
%      throughout the text and in the source tree.
% ---------------------------------------------------------------------------
%  Sharing one counter across environments would normally make cleveref
%  mislabel every cross reference, because it keys off the *counter* name
%  rather than the environment name: \cref to a theorem sharing the
%  `definition' counter prints "Definition". aliascnt gives each environment
%  its own counter name aliased onto a single shared value, so the numbering
%  is shared but cleveref still prints the right word.
\usepackage{aliascnt}

\newcommand{\sharedtheorem}[3][plain]{%
  \theoremstyle{#1}%
  \newaliascnt{#2}{definition}%
  \newtheorem{#2}[#2]{#3}%
  \aliascntresetthe{#2}%
}

\theoremstyle{definition}
\newtheorem{definition}{Definition}[chapter]

\sharedtheorem[definition]{notationx}{Notation}
\sharedtheorem[definition]{example}{Example}
\sharedtheorem[definition]{remark}{Remark}
\sharedtheorem[definition]{convention}{Convention}

\sharedtheorem{axiom}{Axiom}
\sharedtheorem{theorem}{Theorem}
\sharedtheorem{lemma}{Lemma}
\sharedtheorem{proposition}{Proposition}
\sharedtheorem{corollary}{Corollary}
\sharedtheorem{invariant}{Invariant}

% Flat, prefixed counters. These use the upright `definition' style rather
% than the italic `plain' one: their bodies run to several paragraphs and
% contain code, displays and tables, none of which read well in italic.
\theoremstyle{definition}

\newtheorem{restr}{Restriction}
\renewcommand{\therestr}{R\arabic{restr}}

\newtheorem{deviation}{Deviation}
\renewcommand{\thedeviation}{D\arabic{deviation}}

% Styling applied on top of amsthm, so the environments keep their normal
% \begin{restr}\label{..} usage.
\tcolorboxenvironment{restr}{%
  breakable, enhanced, sharp corners=downhill,
  colback=resbg, colframe=dstmblue, boxrule=0pt,
  leftrule=2.5pt, top=6pt, bottom=6pt, left=8pt, right=8pt,
  before skip=10pt, after skip=10pt,
}
\tcolorboxenvironment{deviation}{%
  breakable, enhanced, sharp corners=downhill,
  colback=devbg, colframe=dstmred, boxrule=0pt,
  leftrule=2.5pt, top=6pt, bottom=6pt, left=8pt, right=8pt,
  before skip=10pt, after skip=10pt,
}

% ---------------------------------------------------------------------------
%  Bibliography
% ---------------------------------------------------------------------------
\usepackage[backend=biber,style=numeric,sorting=nyt,maxbibnames=6]{biblatex}
\addbibresource{refs.bib}

% ---------------------------------------------------------------------------
%  Headers, chapter titles, TOC depth
% ---------------------------------------------------------------------------
\usepackage{titlesec}
\titleformat{\chapter}[display]
  {\normalfont\huge\bfseries\color{dstmblue}}
  {\normalsize\MakeUppercase{\chaptertitlename}\ \thechapter}{10pt}{\Huge}
\titlespacing*{\chapter}{0pt}{-20pt}{28pt}

\usepackage{fancyhdr}
\pagestyle{fancy}
\fancyhf{}
\fancyhead[L]{\small\nouppercase{\leftmark}}
\fancyhead[R]{\small\thepage}
\renewcommand{\headrulewidth}{0.4pt}
\fancypagestyle{plain}{\fancyhf{}\fancyfoot[C]{\small\thepage}%
  \renewcommand{\headrulewidth}{0pt}}

\setcounter{tocdepth}{2}
\setcounter{secnumdepth}{3}

% ---------------------------------------------------------------------------
%  Cross references (hyperref must come late, cleveref after it)
% ---------------------------------------------------------------------------
% plainpages=false: the front matter is numbered in roman and the body in
% arabic, so both contain a "page 1" and hyperref would emit a duplicate
% destination without it.
\usepackage[hidelinks,colorlinks=true,plainpages=false,
            linkcolor=dstmblue,citecolor=dstmteal,urlcolor=dstmteal]{hyperref}
\usepackage[capitalise,nameinlink]{cleveref}

\crefname{definition}{Definition}{Definitions}
\crefname{notationx}{Notation}{Notations}
\crefname{example}{Example}{Examples}
\crefname{remark}{Remark}{Remarks}
\crefname{convention}{Convention}{Conventions}
\crefname{axiom}{Axiom}{Axioms}
\crefname{theorem}{Theorem}{Theorems}
\crefname{lemma}{Lemma}{Lemmas}
\crefname{proposition}{Proposition}{Propositions}
\crefname{corollary}{Corollary}{Corollaries}
\crefname{invariant}{Invariant}{Invariants}
\crefname{restr}{Restriction}{Restrictions}
\crefname{deviation}{Deviation}{Deviations}

% ---------------------------------------------------------------------------
%  Glossary
%
%  Loaded after hyperref, as the package requires. No \makeglossaries and no
%  \printglossary: section 1.1 is the glossary, written by hand, and the
%  \term macro plants the \glstarget anchors there. \gls and friends then
%  hyperlink into that section. This also keeps the build a single latexmk
%  invocation with no external makeglossaries step.
% ---------------------------------------------------------------------------
\usepackage[nopostdot,nonumberlist,nogroupskip]{glossaries}
% ============================================================================
%  Glossary entries.
%
%  These declare the *keys* used by \gls{...} throughout the document, and the
%  text each key expands to. The authoritative treatment of every term --- the
%  prose reading plus the formal definition --- lives in section 1.1, and that
%  is where \gls{...} hyperlinks to: the \term macro emits the \glstarget
%  anchor for the corresponding key.
%
%  Consequently no generated glossary is printed. Section 1.1 *is* the
%  glossary; this file exists so that \gls, \glspl, \Gls and \Glspl know what
%  to typeset and where to point.
%
%  To add a term:
%    1. add a \newglossaryentry here;
%    2. add \term{<key>}{<Heading>} at its treatment in 01a-glossary.tex,
%       or \alsoterm{<key>} if it shares a heading with another term.
% ============================================================================

% ---------------------------------------------------------------------------
%  Identity and structure
% ---------------------------------------------------------------------------
\newglossaryentry{process}{
  name={process},
  plural={processes},
  description={One participant in a distributed simulation: a running instance
    of the program. Processes are the entities that own memory, advance along
    the timeline and issue operations}}

\newglossaryentry{rank}{
  name={rank},
  description={An integer in [0,n) naming a process uniquely; the index by
    which every structure in the model is addressed}}

\newglossaryentry{session}{
  name={session},
  description={A maximal interval during which the distance matrix and every
    region size are fixed}}

\newglossaryentry{region}{
  name={region},
  description={The flat, byte-addressable slab of versioned memory owned by
    one process}}

\newglossaryentry{address-space}{
  name={address space},
  plural={address spaces},
  description={The set of byte offsets valid within a given region}}

\newglossaryentry{allocation}{
  name={allocation},
  plural={allocations},
  description={A contiguous run of addresses reserved for one data structure;
    a property of addresses, uniform across all versions}}

\newglossaryentry{pointer}{
  name={pointer},
  description={A (region, offset) pair naming a location; the region
    identifier is also the owning process}}

% ---------------------------------------------------------------------------
%  Time
% ---------------------------------------------------------------------------
\newglossaryentry{logical-time}{
  name={logical time},
  description={Simulation time; the shared timeline along which all processes
    advance, unrelated to wall-clock time}}

\newglossaryentry{origin-time}{
  name={origin time},
  description={The point on the timeline at which the issuing process sits when
    it issues an operation}}

\newglossaryentry{delta}{
  name={delta},
  description={The signed displacement separating when an operation happens
    from when its effect is visible}}

\newglossaryentry{effective-time}{
  name={effective time},
  description={Origin time plus delta; the version an operation actually
    touches}}

\newglossaryentry{local-clock}{
  name={local clock},
  description={The furthest point along the timeline a process has declared
    itself to be at; monotone, and published to neighbours}}

\newglossaryentry{known-time}{
  name={known time},
  description={A process's current lower bound on a neighbour's local clock}}

\newglossaryentry{drift}{
  name={drift},
  description={The clock skew between two processes}}

\newglossaryentry{mvt}{
  name={minimum virtual time},
  description={The oldest version a process still holds}}

% ---------------------------------------------------------------------------
%  Locality
% ---------------------------------------------------------------------------
\newglossaryentry{comm-graph}{
  name={communication graph},
  plural={communication graphs},
  description={The declaration of which pairs of ranks may exchange
    operations, and what each such channel costs in simulated time}}

\newglossaryentry{lag}{
  name={lag},
  plural={lags},
  description={The weight of a channel: the number of simulated timesteps an
    interaction across it takes to become visible. Not a metric}}

\newglossaryentry{path-weight}{
  name={path weight},
  plural={path weights},
  description={The least total lag of a chain of channels joining two ranks;
    a metric, and the bound on how far two processes may drift apart}}

\newglossaryentry{distance-matrix}{
  name={distance matrix},
  plural={distance matrices},
  description={The API-level encoding of a communication graph as an
    n-by-n integer matrix, with -1 marking the absence of a channel}}

\newglossaryentry{reachability}{
  name={reachability},
  description={The relation holding between two processes that the distance matrix
    assigns a distance}}

\newglossaryentry{neighbourhood}{
  name={neighbourhood},
  description={The set of ranks naming the processes reachable from a given process}}

\newglossaryentry{lookahead}{
  name={lookahead},
  description={The guarantee that nothing a process does now can affect a
    neighbour's state sooner than the distance between them}}

\newglossaryentry{heartbeat}{
  name={heartbeat},
  description={A contentless message publishing a process's clock advance to its
    neighbourhood}}

% ---------------------------------------------------------------------------
%  State
% ---------------------------------------------------------------------------
\newglossaryentry{cell}{
  name={cell},
  description={The smallest unit of DSTM state, identified by rank, effective
    time and offset}}

\newglossaryentry{version}{
  name={version},
  description={The whole of one process's region at one effective time; the unit
    the runtime allocates, publishes and reclaims}}

\newglossaryentry{version-status}{
  name={version status},
  plural={version statuses},
  description={A version's position in the one-way progression absent,
    writable, readable, discardable}}

\newglossaryentry{store}{
  name={versioned store},
  description={The total map from versioned cells to bytes; the object the
    memory model is about}}

\newglossaryentry{publication}{
  name={publication},
  description={The promotion of a version to readable, after which its
    contents are immutable}}

% ---------------------------------------------------------------------------
%  Operations
% ---------------------------------------------------------------------------
\newglossaryentry{operation}{
  name={operation},
  description={A one-sided read or write against the versioned store}}

\newglossaryentry{put}{
  name={put},
  description={A write operation, carrying a non-negative delta}}

\newglossaryentry{get}{
  name={get},
  description={A read operation, carrying a non-positive delta}}

\newglossaryentry{footprint}{
  name={footprint},
  description={The set of cells an operation touches, including their
    effective-time coordinate}}

\newglossaryentry{issue}{
  name={issue},
  description={The instant the API call for an operation is made, as distinct
    from its completion}}

\newglossaryentry{completion}{
  name={completion},
  description={The instant an operation's effect on the store, or its result,
    is determined}}

\newglossaryentry{history}{
  name={history},
  plural={histories},
  description={The record of every operation a program issued, totally ordered
    within each process}}

\newglossaryentry{program-order}{
  name={program order},
  description={The total order on the operations issued by a single process}}

% ---------------------------------------------------------------------------
%  Coordination
% ---------------------------------------------------------------------------
\newglossaryentry{quiescence}{
  name={quiescence},
  description={The condition that no operation is in flight and no process will
    issue another}}

\newglossaryentry{barrier}{
  name={barrier},
  description={The three-round ring protocol establishing quiescence}}

\newglossaryentry{termination}{
  name={termination},
  description={Quiescence together with teardown of every process's region}}

% ---------------------------------------------------------------------------
%  Derived structures
% ---------------------------------------------------------------------------
\newglossaryentry{darray}{
  name={distributed array},
  plural={distributed arrays},
  description={A one-dimensional array of fixed-size elements cut into one
    contiguous slice per process}}

\newglossaryentry{partition}{
  name={partition},
  description={The assignment of contiguous global index ranges to processes, named by rank}}

\newglossaryentry{dgraph}{
  name={distributed graph},
  plural={distributed graphs},
  description={A static directed graph carrying versioned state on vertices
    and edges, each edge owned by the owner of its source}}

\newglossaryentry{topology-replica}{
  name={topology replica},
  plural={topology replicas},
  description={The static, unversioned copy of a neighbouring process's adjacency
    structure, exchanged once at construction}}

% ---------------------------------------------------------------------------
%  Idioms
% ---------------------------------------------------------------------------
\newglossaryentry{pull-model}{
  name={pull model},
  description={Producers publish locally at the current time; consumers read a
    neighbour's past}}

\newglossaryentry{push-model}{
  name={push model},
  description={Producers write into a neighbour's future; consumers read their
    own region at the current time}}

\newglossaryentry{mailbox}{
  name={mailbox pattern},
  plural={mailbox patterns},
  description={Indexing a shared structure by (destination, source) so that
    each cell has exactly one writer per effective time}}


% ---------------------------------------------------------------------------
%  Presentation helpers
% ---------------------------------------------------------------------------

% Inline code. C identifiers in this project are long (up to
% dstm_region_version_find_or_create), and TeX will not break inside a
% \texttt run, producing overfull lines. Redefining \_ locally to permit a
% break after each underscore fixes every such identifier at once.
% In math mode the content is wrapped in \mbox so that \ttfamily is legal;
% breaking is lost there, but math-mode occurrences are short identifiers.
\newcommand{\codeus}{\textunderscore\allowbreak}
\newcommand{\codebody}[1]{\small\ttfamily\let\_\codeus #1}
\newcommand{\code}[1]{%
  \ifmmode\mbox{\codebody{#1}}\else{\codebody{#1}}\fi}

% A reference into the implementation, e.g. \srcref{ops.c}{99}.
\newcommand{\srcref}[2]{\texttt{\small #1:#2}}
\newcommand{\srcrefr}[3]{\texttt{\small #1:#2--#3}}

% Inline badge marking a point where the implementation departs from this
% specification; hyperlinks to the register in Appendix B.
\newcommand{\devmark}[1]{%
  \tcbox[on line, boxsep=0pt, left=3pt, right=3pt, top=1pt, bottom=1pt,
         colback=devbg, colframe=dstmred, boxrule=0.5pt, arc=1.5pt]%
        {\scriptsize\bfseries\color{dstmred}\ref{#1}}%
}

% Name of an LTS rule. \textnormal resets the surrounding font first:
% mathpartir sets rule labels in italic, and Latin Modern has no
% small-caps-italic face, so a bare \textsc there triggers a font substitution.
\newcommand{\rn}[1]{\textnormal{\textsc{#1}}}

% Run-in heading for a glossary term. Deliberately not a sectioning command:
% the glossary has ~35 entries and they would swamp the table of contents.
%
%   \term{<key>}{<Heading>}  heading for the term, and the hyperlink target
%                            that every \gls{<key>} in the document points to.
%   \alsoterm{<key>}         an additional target, for a term treated under a
%                            heading belonging to another key.
%
% Exactly one target per key must exist, or hyperref reports a duplicate.
\newcommand{\term}[2]{%
  \par\medskip\noindent
  \glstarget{#1}{}%
  \textbf{\color{dstmblue}\sffamily #2}\par\nopagebreak\noindent%
}
\newcommand{\alsoterm}[1]{\glstarget{#1}{}\ignorespaces}

% Column type: fixed-width, ragged right, for contract tables.
\newcolumntype{L}[1]{>{\raggedright\arraybackslash}p{#1}}

% Table rules a touch lighter than booktabs defaults.
\setlength{\heavyrulewidth}{0.9pt}
\setlength{\lightrulewidth}{0.4pt}

\allowdisplaybreaks
